% ED26-SCI-Mahajan2014-001; dimensional constraints and scaling inference.
\chapter{Dimensional and Scaling Inference}
\label{sec:p15-dimensional-and-scaling-inference}

Dimensional reasoning constrains a proposed relation by requiring consistent
behavior under changes of units. Scaling reasoning asks how a quantity changes
when selected features of the problem change. The first concerns representation;
the second concerns a family of situations. Neither supplies an omitted
physical premise merely by eliminating inadmissible formulas.

This treatment reconstructs Mahajan's 2014 Chapter 4, Sections 5.1--5.3 and
the estimation methods in Sections 1.3--1.5 \cite{Mahajan2014Insight}.
The linear-algebra and scaling theorems below state their domains explicitly.
Physical laws used in examples are declared model premises. Numerical source
illustrations retain their historical standing; they are not new measurements
or current engineering specifications.

\section{Quantities, units and relevant comparisons}
\label{sec:p15-dimensional-quantities}

\begin{hbdefinition}{def:p15-dimensional-coordinate}{Dimensions and numerical values}
Choose \(k\) multiplicative base-unit types. Assign a quantity \(q_i\) a
dimension vector \(d_i\in\mathbb R^k\), whose coordinates are its powers
of those types. Its numerical value is \(x_i\).
For example, with length, mass and time in that order,
\[
 [v]=(1,0,-1),\quad [a]=(1,0,-2),\quad
 [F]=(1,1,-2),\quad [\rho]=(-3,1,0).
\]
Brackets here mean the dimension vector, not a physical value or an interval.
Products add dimension vectors; a real power multiplies the vector by that
power on a domain where the power is defined.

If the new base units are \(e^{-s_1},\ldots,e^{-s_k}\) times the old ones,
the new numerical values are
\[
 x_i'=x_i e^{d_i^{\mathsf T}s}.
\]
Positive quantities and multiplicative unit changes are the initial domain
of the logarithmic formulation below.
\end{hbdefinition}

Lengths can be expressed in metres or millimetres without changing which
physical length is larger. A numerical value such as \(540\) does not by
itself establish that a measured rainfall is large. Mahajan compares
\(540\,\mathrm{mm/year}\) with both \(1000\,\mathrm{mm/year}\) and
\(350\,\mathrm{mm/year}\), obtaining ratios \(0.54\) and about \(1.54\).
The different conclusions answer different reference questions.
Comparing the same rainfall rate with an orbital speed is dimensionally
possible but does not answer a useful wetness question
\cite[pp.~137--138]{Mahajan2014Insight}. Equal dimensions are a compatibility
check, not proof of semantic relevance.

Addition requires compatible dimensions and meanings: an energy can be added
to another energy, not to a bare speed. A logarithm or exponential ordinarily
takes a dimensionless numerical ratio, such as \(\log(x/x_0)\).
A named coefficient can carry dimensions; dropping it may change the
dimensional type of an equality.

Affine temperature scales require separate care. A Celsius reading and an
absolute temperature do not differ only by multiplication. The multiplicative
dimension theorem applies to an absolute scale, or to temperature differences
with their stated conversion. A dimensionless angle also retains a numerical
convention: radians and degrees differ by a constant even though both have
zero physical dimension.

\section{Dimensionless products and the rank condition}
\label{sec:p15-dimensional-groups}

\begin{hbdefinition}{def:p15-dimensional-matrix}{Dimension matrix and product}
For \(n\) positive quantities, form the \(k\)-by-\(n\) matrix
\(D=(d_1\ \cdots\ d_n)\). For \(a\in\mathbb R^n\), define
\[
 x^a=\prod_{i=1}^n x_i^{a_i}
       =\exp\!\left(\sum_i a_i\log x_i\right).
\]
A product is \emph{dimensionless} when \(Da=0\).
Independent dimensionless products mean linearly independent exponent vectors.
The rank of \(D\) is the dimension of its column space; its kernel is
the set of exponent vectors sent to zero.
\end{hbdefinition}

\begin{hbtheorem}{thm:p15-dimensional-rank-nullity}{proved=annex:p15.dimensional.rank-nullity}{Counting and generating the independent products}
If \(D\) has rank \(r\), its dimensionless exponent vectors form a vector
space of dimension \(n-r\). A basis \(b_1,\ldots,b_{n-r}\) gives products
\(\pi_j=x^{b_j}\), and every dimensionless monomial has the form
\(\prod_j\pi_j^{c_j}\). Invertible changes of dimension coordinates do not
change this space or its dimension.
\end{hbtheorem}

This is the precise counting argument behind the source's Buckingham-Pi
discussion \cite[5.1.2]{Mahajan2014Insight}. Count independent dimensions,
not the number of dimension symbols that happen to occur.

For the three quantities weight \(W\), mass \(m\) and acceleration \(g\),
the matrix in the order \((\mathrm L,\mathrm M,\mathrm T)\) is
\[
 D=\begin{pmatrix}1&0&1\\1&1&0\\-2&0&-2\end{pmatrix}.
\]
It has rank two: its third row is \(-2\) times the first. Solving \(Da=0\)
gives \(a=t(1,-1,-1)\), so \(W/(mg)\) is one independent product.
The appearance of three dimension names does not give rank three.

For acceleration \(a\), speed \(v\) and radius \(r\),
\[
 D=\begin{pmatrix}1&1&1\\-2&-1&0\end{pmatrix},
 \qquad \ker D=\{t(1,-2,1):t\in\mathbb R\}.
\]
The product \(ar/v^2\) is therefore a basis product. Its value is not
determined by the kernel computation.

\begin{hbmethod}{def:p15-dimensional-elimination}{Finding a useful dimensional basis}
\begin{enumerate}
\item State the output, candidate inputs, units, domain and fixed parameters.
Include dimensional constants actually used by the model.
\item Put dimension vectors in columns and reduce the matrix to row-echelon
form. Keep its rank and its free exponent variables.
\item Assign the free variables one at a time to form a nullspace basis.
Check every proposed product by multiplication with \(D\).
\item Choose an equivalent basis convenient for the question. When possible,
put the output in one product and keep it out of the others.
\item State the remaining arbitrary function or constant. Supply additional
model, experimental or boundary information before assigning it a value.
\end{enumerate}
\end{hbmethod}

The positivity restriction is substantive. Logs and arbitrary real powers
do not represent zero or negative quantities. On a fixed nonzero sign region
one can use absolute values and retain signs as extra invariants. Zero
patterns form separate cases. Dividing a proposed relation by an expression
that can vanish can discard solutions.

\section{What the Pi representation actually establishes}
\label{sec:p15-dimensional-pi}

\begin{hbtheorem}{thm:p15-dimensional-pi-factorization}{proved=annex:p15.dimensional.pi-factorization}{Complete invariants on the positive domain}
Let \(B\) have as rows a basis of \(\ker D\), and put
\(\pi(x)=\exp(B\log x)\), with componentwise logarithms and exponentials.
A function \(F:(0,\infty)^n\to S\), into any set \(S\), is invariant under
all unit changes
\[
 x\longmapsto x\odot\exp(D^{\mathsf T}s)
\]
if and only if it has the form \(F(x)=\Psi(\pi(x))\).
Here \(\odot\) denotes componentwise multiplication.
The same statement holds for membership in an invariant set of allowed
configurations.
\end{hbtheorem}

\begin{hbtheorem}{thm:p15-dimensional-covariant-output}{proved=annex:p15.dimensional.covariant-output}{A dimensionally consistent output function}
Suppose a numerical output \(f(x)\) has dimension vector \(w\) and satisfies
\[
 f(x\odot\exp(D^{\mathsf T}s))=e^{w^{\mathsf T}s}f(x)
 \quad\text{for every }x>0,\ s\in\mathbb R^k.
\]
If \(Da=w\) has a solution, then
\[
 f(x)=x^a\Psi(\pi(x)).
\]
If it has no solution, the only real-valued function satisfying the displayed
covariance for all unit changes is the zero function.
\end{hbtheorem}

A missing solution can therefore reveal that a proposed nonzero output needs
another dimensional input or a changed specification. It does not identify
which physical quantity was omitted.

Even a one-product relation need not select one constant. For example,
\((\pi-1)(\pi-2)=0\) allows two dimensionless branches, and
\(1\leq\pi\leq2\) allows a range. A universal single constant follows for
a single-valued covariant output with no remaining dimensionless input
and a fixed model family. Branch selection, dimensionless shape, boundary
conditions and constitutive parameters must not be suppressed.

\section{Scaling laws and the information in a log plot}
\label{sec:p15-dimensional-scaling}

\begin{hbdefinition}{def:p15-dimensional-scaling-law}{A declared scaling family}
Let \(x,y>0\). Within a specified family, \(y=Kx^\alpha\), with constant
\(K>0\), means
\[
 \frac{y_2}{y_1}=\left(\frac{x_2}{x_1}\right)^\alpha.
\]
The held-fixed conditions are part of this claim. The coefficient \(K\)
need not be dimensionless. In this chapter \(\propto\) denotes such a
constant-ratio relation; a rough estimate is explicitly labelled approximate.
\end{hbdefinition}

\begin{hbtheorem}{thm:p15-dimensional-scaling-calculus}{proved=annex:p15.dimensional.scaling-calculus}{Composition and branching of exponents}
Take positive variables and coefficients, and real exponents, in the
following finite products.
If \(u=Ax^\alpha\) and \(y=Bu^\beta\), then
\(y=BA^\beta x^{\alpha\beta}\).
If \(y=K\prod_j u_j^{p_j}\), where
\(u_j=A_j\prod_i x_i^{a_{ji}}\), then
\[
 y=K\prod_j A_j^{p_j}\prod_i x_i^{\sum_jp_j a_{ji}}.
\]
Thus exponent factors multiply along a path and contributions add where
multiplicative paths meet. These rules do not replace addition of quantities
by addition of exponents.
\end{hbtheorem}

Mahajan's dinner example doubles \(250\,\mathrm g\) for four diners to
\(500\,\mathrm g\) for eight under unchanged consumption per diner.
His historical gas-station estimate scales five stations for 20,000 people
to 75,000 for a stipulated population of \(3\,10^8\). Its evidence concerns
the station-per-person assumption and the input estimates, not a theorem
that all towns have the same ratio \cite[4.1]{Mahajan2014Insight}.
The source compares that rough estimate with a reported 2008 count of
116855; agreement at that precision does not establish population
proportionality in every locality.

\begin{hbtheorem}{thm:p15-dimensional-log-slope}{proved=annex:p15.dimensional.log-slope}{Logarithmic slopes and their limits}
Fix positive references \(x_0,y_0\). A power law is a straight line on
the coordinates \(\xi=\log(x/x_0)\), \(\eta=\log(y/y_0)\), with slope
\(\alpha\). For differentiable positive \(f\), the local slope is
\(xf'(x)/f(x)\). If this slope is constantly \(\alpha\) on an interval,
then \(f(x)=Kx^\alpha\) there.
A single doubling relation does not establish a power law.
\end{hbtheorem}

For example,
\[
 f(x)=x^\alpha\exp\{\varepsilon\sin(2\pi\log_2x)\},\qquad\varepsilon\ne0,
\]
satisfies \(f(2x)=2^\alpha f(x)\) but has a varying log slope.
A straight segment observed over a finite range is evidence about that
range, not an exact scaling law at every scale.

The source's inverse-square gravity plot has log slope \(-2\).
Its Zipf illustration uses word frequency proportional to inverse rank,
and its population illustration uses a similar rank estimate.
These are different model families: a rank-frequency regularity is not
deduced from dimensional consistency \cite[4.5]{Mahajan2014Insight}.
The source's illustrative word table gives frequencies 7.0, 3.5 and 2.9
percent at ranks one, two and three. Its state-population example uses a
40-million reference and a putative rank near 50, giving roughly one
million, compared with a reported 917000. Actual rank and the applicability
of the rank law must be established; small land area does not determine
population rank.

\section{Geometric and mechanical source cases}
\label{sec:p15-dimensional-mechanical}

\begin{hbtheorem}{thm:p15-dimensional-mechanical-models}{proved=annex:p15.dimensional.mechanical-models}{Consequences of explicit geometric and motion models}
The following conclusions hold under their stated models.
\begin{enumerate}
\item Uniform scaling of a Borel planar set of positive finite area by \(s>0\) multiplies its area by
\(s^2\). Halving the area of a similar figure multiplies its linear dimensions
by \(1/\sqrt2\).
\item Uniform circular motion of radius \(r>0\) and angular frequency \(\omega>0\)
has speed \(v=r\omega\), inward acceleration \(v^2/r\), and period
\(2\pi/\omega\).
\item A circular test-particle orbit under acceleration \(GM/r^2\), with
fixed \(GM>0\), has period \(2\pi\sqrt{r^3/(GM)}\).
\item A projectile with no drag, constant downward acceleration \(g>0\),
launch speed \(v>0\), angle \(0<\theta<\pi/2\), and equal launch/landing
height has range \(v^2\sin(2\theta)/g\).
\item The linear oscillator \(\ddot x+(k/m)x=0\), \(k,m>0\), has
period \(2\pi\sqrt{m/k}\) for a nonzero oscillation.
\end{enumerate}
\end{hbtheorem}

The source's circle has diameter \(\sqrt5\,\mathrm{cm}\), so its half-area
similar circle has diameter \(\sqrt{5/2}\,\mathrm{cm}\), not half the original
diameter \cite[4.2.1]{Mahajan2014Insight}.
All planar shapes obey the \(s^2\) relation under uniform scaling; changing
shape at the same time need not preserve the proportionality coefficient.

The circular-orbit model gives the exponent \(3/2\), and a radius ratio
\(1.52\) gives a period ratio \(1.52^{3/2}\approx1.874\).
Applying this to an eccentric orbit requires the appropriate orbital model
and semimajor axis, rather than an arbitrary instantaneous distance.
Test-particle mass cancels under the declared fixed-central-field model;
a two-body model has a different mass parameter \cite[4.2.2]{Mahajan2014Insight}.

Surface gravity illustrates the importance of what is fixed. Under
\(g=GM/R^2\), varying \(R\) at fixed \(M\) gives exponent \(-2\).
With uniform-sphere mass \(M=(4\pi/3)\rho R^3\), varying \(R\) at fixed
\(\rho\) instead gives exponent \(+1\):
\[
 g=\frac{4\pi}3G\rho R.
\]
These are compatible conditional claims, not rival values of one unspecified
exponent \cite[4.2.4]{Mahajan2014Insight}.

For the source's radiation illustration, isotropic luminosity held fixed
gives flux \(F=L/(4\pi r^2)\). A declared radiative balance
\(\epsilon\sigma T^4=(1-A)F/4\), with unchanged emissivity \(\epsilon\)
and albedo \(A\), gives \(T\propto r^{-1/2}\).
Changes in these factors, internal heating or atmospheric assumptions defeat
an unrestricted planet-to-planet extrapolation.
Using the source's \(F\approx0.8\,\mathrm{W/m^2}\) and
\(\sigma\approx5.7\,10^{-8}\,\mathrm{W\,m^{-2}K^{-4}}\), with
\(A=0,\epsilon=1\), gives about \(43.3\,\mathrm K\).
The alternative comparison \(293/\sqrt{40}\) gives about \(46\,\mathrm K\)
and uses additional reference-planet assumptions
\cite[4.2.1]{Mahajan2014Insight}.

\section{A period scale, its constant and an omitted amplitude}
\label{sec:p15-dimensional-pendulum}

For pendulum period \(T\), length \(l\), gravity \(g\) and bob mass \(m\),
the dimension matrix has four columns and rank three. The only independent
product can be chosen as \(T\sqrt{g/l}\); mass cannot enter a product when
it is the only quantity carrying a mass dimension.
A single-valued model with only \(l,g,m\) therefore has
\[
 T=C\sqrt{l/g}.
\]
This establishes neither the constant nor amplitude independence.
An amplitude \(\theta_0\) is an additional dimensionless input, allowing
\(T=\sqrt{l/g}\,\Psi(\theta_0)\).

For a point bob on a massless rigid-length support in uniform gravity,
linearizing \(\sin\theta\) by \(\theta\) gives
\(\ddot\theta+(g/l)\theta=0\).
The oscillator result then gives \(C=2\pi\) in that linear model.
The full pendulum equation retains amplitude dependence. Small amplitude
is an approximation premise, not a conclusion of the dimensions
\cite[5.2.1]{Mahajan2014Insight}.

Mahajan reports a key-ring estimate with \(l\approx0.65\,\mathrm m\)
and ten periods taking \(15.97\,\mathrm s\); using \(g=9.8\,\mathrm{m/s^2}\)
gives \(T\sqrt{g/l}\approx6.201\).
This is a calibration illustration, not a proof that the exact value is
\(2\pi\). Feet/seconds and metres/seconds give the same dimensionless
quantity when conversions are consistent. In the curved-track example,
\(v=60\,\mathrm{m/s}\), \(r=2000\,\mathrm m\) give \(a=1.8\,\mathrm{m/s^2}\);
\(a/g\approx0.18\) corresponds to \(\tan\theta\approx0.18\).
The source's walking model similarly uses \(v^2/(gl)\), a Froude-number
convention, while its transition constant requires biomechanical evidence
\cite[5.1.1--5.2.1]{Mahajan2014Insight}.

Adding an initial downward speed to a falling-body problem creates another
product:
\[
 \frac{v}{\sqrt{gh}}=\Psi\!\left(\frac{v_0}{\sqrt{gh}}\right).
\]
The constant-acceleration or energy model supplies
\(\Psi(z)=\sqrt{z^2+2}\); dimensions alone do not.
The cases \(h=0\), \(g=0\), or a changed force law require separate domains
\cite[5.3]{Mahajan2014Insight}.

\section{Similarity, drag and fixed-regime conditions}
\label{sec:p15-dimensional-drag}

For a fixed shape, with drag force \(F_{\rm d}\), density \(\rho\),
speed \(v\), area \(A\) and kinematic viscosity \(\nu\),
dimensional elimination gives two independent products:
\[
 c_{\rm d}=\frac{2F_{\rm d}}{\rho v^2 A},\qquad
 \mathrm{Re}=\frac{vD_{\rm b}}{\nu}.
\]
Here \(D_{\rm b}\) is a declared characteristic length, proportional to
\(\sqrt A\) for the fixed shape. The general single-valued form is
\(c_{\rm d}=\Psi(\mathrm{Re})\). Shape parameters, compressibility,
surface condition or other governing quantities add further dimensionless
arguments when relevant \cite[5.3.2]{Mahajan2014Insight}.
Omitting viscosity from the list does not prove the drag constant is zero;
it leaves a physically incomplete candidate model.

For quadratic drag and terminal balance \(mg=\frac12c_{\rm d}\rho A v^2\),
\[
 v=\sqrt{\frac{2mg}{c_{\rm d}\rho A}}.
\]
For similar cones cut from paper of fixed areal density, \(m\) and \(A\)
both scale as length squared. Their terminal speeds are therefore equal
\emph{if} the drag coefficients and remaining conditions agree.
The coefficients cannot be cancelled merely because the shapes agree when
their Reynolds numbers differ \cite[4.3.1]{Mahajan2014Insight}.

The source reports an informal falling-cone comparison, and estimates
landing times within about 0.1 seconds of one another after a roughly
2-metre, 2-second fall. It does not report a time series establishing the
terminal regime. It then estimates
\(\mathrm{Re}\) near \(3500\) and \(7000\) for two sizes at about
\(1\,\mathrm{m/s}\), with \(\nu\approx1.5\,10^{-5}\,\mathrm{m^2/s}\).
Using paper areal density \(0.08\,\mathrm{kg/m^2}\),
\(A=(3/4)A_{\rm paper}\), and \(\rho=1.2\,\mathrm{kg/m^3}\) gives
\(c_{\rm d}\approx1.8\). These are source-specific approximate observations
and parameter estimates, not a universal cone coefficient.

The sphere plot at printed page 164 shows a roughly flat region followed
by a marked drag-coefficient decrease near Reynolds number \(3\,10^5\).
It demonstrates why a plateau-based scaling relation must carry its regime.
A fixed coefficient inferred from a few trials is not established for all
speeds or sizes. Equal fall times from a finite height also require attention
to acceleration toward terminal speed \cite[4.3.1, 5.3.2]{Mahajan2014Insight}.

\section{Competing powers and flight comparisons}
\label{sec:p15-dimensional-optimization}

\begin{hbtheorem}{thm:p15-dimensional-competing-powers}{proved=annex:p15.dimensional.competing-powers}{The minimum of two opposing powers}
For \(A,B,p,q>0\), the function
\(E(v)=Av^{-p}+Bv^q\), \(v>0\), has a unique minimum at
\[
 v_*=\left(\frac{pA}{qB}\right)^{1/(p+q)}.
\]
At that point \(pAv_*^{-p}=qBv_*^q\).
The two contributions are equal only when \(p=q\).
For \(E(v)=A/v^2+Bv^2\), \(v_E=(A/B)^{1/4}\) and
\(E(v_E)=2\sqrt{AB}\). If power is \(P(v)=vE(v)\), then
\(v_P=3^{-1/4}v_E\).
\end{hbtheorem}

For energy per unit flight distance, Mahajan uses opposing model contributions with
\[
 A=\frac{(mg)^2}{\rho L^2},\qquad B=C\rho L^2,
\]
where \(L\) is span and \(C\) is a drag coefficient referenced to \(L^2\).
These describe induced and profile energy costs in a simplified model;
they do not assert equality of vertical lift and horizontal drag.
At the model's energy optimum, the total energy per distance is
\(2\sqrt C\,mg\), and each contribution is \(\sqrt C\,mg\)
\cite[4.6]{Mahajan2014Insight}.

The source defines three area conventions, related by
\[
 CL^2=\tfrac12 C' A_{\rm wing}=\tfrac12c_{\rm d}A_{\rm cs}.
\]
Thus a quoted coefficient cannot be transferred without its reference area
and factor-of-two convention. Using the source's
\(C'\approx0.022\), wing chord \(10\,\mathrm m\) and span \(60\,\mathrm m\)
gives \(C\approx1/600\).
Its order estimate discards a factor of order two and uses
\[
 d\sim\frac{\beta\mathcal E_{\rm fuel}}{\sqrt C\,g},
\]
with fuel fraction \(\beta\) and usable energy per fuel mass
\(\mathcal E_{\rm fuel}\). The inputs \(0.4\), \(10^7\,\mathrm{J/kg}\)
and \(g\approx10\,\mathrm{m/s^2}\) give about \(9.8\,10^6\,\mathrm m\).
The exact constant within the two-term surrogate and the source's rough
range estimate have different precision commitments.
Changing mass during fuel consumption is another modeling issue.

At fixed shape, fuel fraction, usable energy density, coefficient regime
and gravity, this approximate range is independent of size.
For geometrically similar masses, \(m\propto L^3\), the same model gives
\(v_E\propto m^{1/6}\rho^{-1/2}\), with other dimensionless parameters fixed.
These conditions also delimit the source's plane/bird comparison.

The source's separate car/plane estimate compares drag energy over the same
distance. Under quadratic drag, its ratio includes the drag-coefficient ratio
as well as density, area and squared speed ratios; converting energy to fuel
also requires the efficiency and fuel-energy conventions. With comparable
coefficients and conversions, density ratio \(1/3\), area ratio \(10\) and speed
ratio \(10\) give a rough vehicle-level factor \(300\).
Dividing by a stipulated \(300\) passengers changes the per-person comparison.
Its quoted 747-400 figures, \(216840\,\mathrm L\), \(13450\,\mathrm{km}\)
and 416 passengers, imply about \(3.88\,\mathrm L\) per
100 passenger-kilometres. Occupancy, efficiency, coefficients and the
distance convention must be specified; these historical inputs are not a
current travel comparison \cite[4.3.2]{Mahajan2014Insight}.

\section{A scaling argument in finite probability}
\label{sec:p15-dimensional-birthday}

\begin{hbtheorem}{thm:p15-dimensional-birthday}{proved=annex:p15.dimensional.birthday}{Birthday collisions and the limits of pair counting}
Let integer \(n\geq1\) independent birthdays be uniform on integer \(M\geq2\) categories.
For \(n\leq M\), the probability of no repeated category is
\[
 P_0=\prod_{j=0}^{n-1}(1-j/M).
\]
For \(n>M\), it is zero. For \(n\leq M\), set
\(S_1=n(n-1)/(2M)\), \(a=(n-1)/M<1\). Then
\[
 -S_1-\frac{\sum_{j=0}^{n-1}(j/M)^2}{2(1-a)}
 \leq\log P_0\leq-S_1.
\]
For an integer threshold of \(k\geq2\) equal birthdays, the exact noncollision
probability is
\[
 P(\max_b N_b<k)=\frac{n!}{M^n}[z^n]
       \left(\sum_{j=0}^{k-1}\frac{z^j}{j!}\right)^M.
\]
Here \(N_b\) is the occupancy of category \(b\), and \([z^n]\) means
the coefficient of \(z^n\).
\end{hbtheorem}

Pair counting makes the scale \(n^2/M\) visible. It does not make all
pair-collision events jointly independent: if persons 1 and 2 and persons
2 and 3 share a birthday, persons 1 and 3 do too.
The exact product and its log bounds justify the approximation
\(P_0\approx e^{-n(n-1)/(2M)}\) in a regime with small remainder
\cite[4.4]{Mahajan2014Insight}. For \(M=365\), the exact collision
probabilities at \(n=22,23\) are approximately \(0.475695\) and \(0.507297\).
Thus the first 50-percent crossing is 23 under this model, not under every
real birthday distribution.

For triples, the exact occupancy expression gives probabilities about
\(0.499455\) and \(0.511065\) at \(n=87,88\).
The heuristic scale \(n^k/M^{k-1}\) for a \(k\)-fold collision comes from
counting \(k\)-subsets: their expected count is
\(\binom nk/M^{k-1}\). Expected count is not collision probability,
though it gives a union bound. This distinction is essential when extending
the source's pair argument to its further questions.

\section{Several products and the choice of physical premises}
\label{sec:p15-dimensional-multiple-groups}

The source's idealized point-explosion example selects energy \(E\), time
\(t\), radius \(R\) and ambient density \(\rho\). Their dimension matrix
has rank three and one product \(Et^2/(\rho R^5)\).
For an energy-conserving similarity regime with fixed \(\rho,E\), a
single selected branch gives \(R\propto t^{2/5}\).
The source's historical radius/time data illustrate that slope; calibration
of the dimensionless coefficient and the regime assumptions are additional
information. The dimensional calculation alone is not an exact energy
determination \cite[5.2.2]{Mahajan2014Insight}.

The four reported pairs \((t\text{ in ms},R\text{ in m})\) are
\((3.26,59.0)\), \((4.61,67.3)\), \((15.0,106.5)\) and
\((62.0,185.0)\). With time converted to seconds,
\(R^5/t^2\) is approximately \(6.73,6.50,6.09,5.64\) times
\(10^{13}\,\mathrm{m^5/s^2}\). Multiplication by a stipulated
\(1\,\mathrm{kg/m^3}\) gives an energy scale in joules, still multiplied
by the undetermined dimensionless coefficient. Agreement of the four scales
supports the proposed time exponent over these observations; it does not
measure that coefficient or establish the complete similarity model.

The light-deflection example illustrates an omitted quantity. An angle
\(\theta\), gravitational parameter \(GM\), and distance \(r\) allow no
dimensionless combination of \(GM\) and \(r\) alone.
Adding speed \(c\) gives
\[
 z=\frac{GM}{rc^2},\qquad \theta=\Psi(z)
\]
for the specified single-valued model. Dimensions leave \(\Psi\) free.
Oddness and monotonicity do not force \(\Psi(z)\propto z\):
both \(z\) and \(z^3\) are odd and increasing.
A differentiable weak-field law with \(\Psi(0)=0\) and
\(\Psi'(0)\ne0\) instead has a first-order term \(\Psi'(0)z\).
The derivative and its value come from extra theory
\cite[5.3.1]{Mahajan2014Insight}.

Mahajan compares coefficient 2 from a Newtonian trajectory approximation
with coefficient 4 from weak-field general relativity,
and estimates a solar-to-terrestrial ratio of order \(3000\) using
mass and radius ratios. That comparison is conditional on the weak-field
model and the chosen ray geometry. A diffraction estimate
\(\lambda/D_{\rm aperture}\) concerns a particular resolution model; it
does not prove that every astrometric technique is incapable of detecting a
smaller positional change. Dimensional constraints do not themselves
establish either a theory's historical confirmation or an observational
impossibility.

\section{Finite estimates and their uncertainty status}
\label{sec:p15-dimensional-estimates}

\begin{hbtheorem}{thm:p15-dimensional-product-bounds}{proved=annex:p15.dimensional.product-bounds}{A product estimate with explicit input ranges}
Let \(Q=K\prod_i x_i^{a_i}\), with \(K>0\) and
\(0<l_i\leq x_i\leq u_i\).
Its exact minimum and maximum over that rectangular domain are obtained
by choosing \(l_i,u_i\) according to the sign of \(a_i\).
If \(x_i=\widehat x_i e^{\delta_i}\), \(|\delta_i|\leq\eta_i\), then
\[
 \left|\log\frac{Q}{K\prod_i\widehat x_i^{a_i}}\right|
 \leq\sum_i|a_i|\eta_i.
\]
These are deterministic range bounds. They are not confidence or posterior
intervals without an additional probability model.
\end{hbtheorem}

Mahajan's tree notation places the target at the root, subquantities below
it and exponent labels on edges. A paper-thickness node, for example, divides
ream thickness by sheets per ream. With \(15\,\mathrm{cm}\),
\(6\,\mathrm{cm}\), \(5\,\mathrm{cm/ream}\) and \(500\) sheets per ream,
the bill-volume model gives \(0.9\,\mathrm{cm^3}\), reported at order-one
precision. Multiplication of the leaf values is justified by the model,
not by the shape of the diagram \cite[1.3]{Mahajan2014Insight}.

For a train line, the source multiplies 150 people per car, 20 cars per train
and six trains per hour to obtain 18,000 people per hour.
A deliberately different illustrative range model,
\[
 100\leq p\leq150,\quad 8\leq c\leq12,\quad4\leq h\leq6,
\]
gives \(3200\leq pch\leq10800\) people per hour.
This range is conditional on feasible combinations and actual occupancy.
If capacity constraints prevent simultaneous endpoint choices, the
rectangular upper bound need not be attained.
Such a range is also conditional on the model omitting no relevant input.
\Cref{sec:p15-w02b-p15-073} works a lane-capacity estimate with input
ranges and exhibits an omitted quantity, vehicle length divided by speed,
that in slow traffic puts the target outside the stated range and makes the
estimate wrong by more than an order of magnitude.

The source's oil estimate has the structure
\[
 Q=N_{\rm people}\frac{N_{\rm cars}}{N_{\rm people}}
 \frac{\text{distance}}{\text{car year}}
 \left(\frac{\text{distance}}{\text{gallon}}\right)^{-1}
 \left(\frac{\text{gallons}}{\text{barrel}}\right)^{-1}
 \frac{\text{all use}}{\text{car use}}
 \frac{\text{imports}}{\text{all use}}.
\]
Its historical values \(3\,10^8,1,20000,30,60,2,1/2\) give about
\(3.33\,10^9\) barrels per year, rounded by the source to \(3\,10^9\).
Its comparison with a reported 2010 rate is a historical illustration.
The reported \(9.163\) million barrels per day corresponds, on a
365-day conversion, to about \(3.34\) billion barrels per year.
The unit conversion used in the estimate is roughly four \emph{litres per
US gallon}, not four gallons per litre \cite[1.4]{Mahajan2014Insight}.
An oil barrel, a roadside container and a fuel price ratio need not refer
to the same volume or product composition.

A second estimate can expose a wrong assumption, but agreement between
methods sharing an input error does not establish accuracy.
The source's reversed-addition and alternative-barrel examples illustrate
the value of changing the route of calculation \cite[1.5]{Mahajan2014Insight}.
For a threshold conclusion \(Q>q_*\), the range model establishes it only
when its justified lower bound exceeds \(q_*\). Point-estimate agreement
alone is insufficient.

\section{Additional source tasks and their inference boundaries}
\label{sec:p15-dimensional-source-tasks}

The selected source includes further questions. The following groups retain
their reasoning tasks, distinguishing calculations from experiments or
external observations that the source asks a reader to undertake.
They do not report those proposed measurements as completed.

\begin{itemize}
\item Source Problems 4.1--4.9 ask for per-person rate comparisons,
a triangle's half-area cut, exponent composition, radiative assumptions,
circular-orbit periods, pendulum mass dependence and projectile prefactors.
The relevant denominators, similarity and force premises are the ones
stated in the preceding examples.
\item Problems 4.10--4.14 connect an orbital-period ratio to distance,
an angular diameter to a radius, density assumptions to surface gravity,
the interior/exterior field of a uniform sphere, and fall height to rotation
angle. Their conditional formulas are
\(r/r_0=(T/T_0)^{2/3}\), \(R\approx r\theta/2\) for small angular diameter,
\(g\propto\rho R\), interior \(g\propto r\), exterior \(g\propto r^{-2}\),
and rotation angle proportional to \(\sqrt h\) at fixed angular speed.
\item Problems 4.15--4.22 concern tape mass in scaled cones, larger-cone
experiments, a four-cone stack, area estimates, fluid-density comparisons,
raindrops, fuel cost and cycling energy.
At fixed quadratic-drag coefficient a fourfold mass increase with unchanged
area doubles terminal speed; a similar solid drop gives \(v\propto r^{1/2}\).
Tape length ordinarily scales as \(r\), paper mass as \(r^2\).
A density comparison based on equal drag power uses \(\rho A v^3\);
a fuel-cost estimate additionally needs distance, occupancy and price.
\item Problems 4.23--4.28 and 4.39 ask for birthday probabilities,
higher-order collisions, bubble sort, merge sort, school multiplication
and Karatsuba scaling. The exact occupancy formula above supplies the
collision problem; the recurrence result below supplies the declared
algorithmic cost models.
\item Problems 4.29--4.38 and 4.40 concern log plots, a population-based
budget estimate, energy versus power optima, lift coefficients, geometric
flight scaling, air density, a migrating-bird comparison, thermal resistance
and a housefly's power. The geometric model gives
\(v_E\propto m^{1/6}\rho^{-1/2}\).
With hover-power model \(P\propto(mg)^{3/2}/(\sqrt\rho L)\),
\(P/m\propto L^{1/2}\) for similar constant-density bodies.
A thermal-resistance comparison requires thickness divided by conductivity
and area. Population, prices, observed speeds and material data remain inputs,
not dimensional conclusions.
\item Problems 5.7--5.12 ask for group counts and changes from speed to
angular frequency or period. Circular motion gives \(a/(r\omega^2)=1\)
and \(aT^2/r=4\pi^2\); dropped-body speed has product \(v/\sqrt{gh}\).
An oscillator with \(T,k,m,x_0,g\) has two independent products, for example
\(T\sqrt{k/m}\) and \(kx_0/(mg)\).
\item Problems 5.13--5.18 ask for pendulum and shallow-wave measurements,
an ideal or nonideal spring, thrown-body speed and rolling motion.
Shallow-wave dimensions permit \(v=C\sqrt{gh}\); the constant and the
long-wavelength regime need physical information.
A quartic potential \(V(x)=Cx^4\) gives period proportional to
\(\sqrt{m/C}/x_0\) for positive amplitude. Rolling under the declared
no-slip force/torque model gives \(a=g\sin\theta/(1+I/(mr^2))\).
\item Problems 5.19--5.24 concern \(GM\), interpretation of
\(GM/(rc^2)\), observational baselines, alternative gravitational models,
large cones and Reynolds-number estimates. The gravity parameter has
dimension \(\mathrm L^3\mathrm T^{-2}\).
Alternative-theory coefficients require that theory; neither dimensions
nor an odd-function constraint supplies them.
\item Source Problems 1.5--1.7 ask for an estimate tree, metric re-expression
and aircraft fuel estimates. All retain the distinction between a conditional
product calculation and an independently checked empirical input.
\end{itemize}
\cite[1.3--1.5, 4.1--4.7, 5.1--5.3]{Mahajan2014Insight}

Some numerical and representational details of these tasks make their
assumptions especially visible:
\begin{itemize}
\item Problem 4.1 supplies historical annual counts of about 14000 and 640;
the per-person ratio is \((14000/640)(N_{\rm UK}/N_{\rm US})\), so population
denominators cannot be omitted. Problem 4.2's horizontal equal-area cut of
an equilateral triangle has length equal to the side divided by \(\sqrt2\).
\item Problem 4.10 supplies a near-surface orbital period of about
1.5 hours. A chosen lunar period must be converted to the same time unit
before applying the \(2/3\) power. Problems 4.11--4.12 additionally need an
observed angular diameter and an explicit density comparison.
For Problem 4.14, changing a half-turn to a full turn at the same angular
speed requires twice the fall time and four times the height in the
zero-initial-vertical-speed model.
\item In Problems 4.15--4.17, doubling cone radius requires four times the
paper mass but only twice the length of an ordinary seam. Added tape must
be included in the mass scaling. Stacking four cones on one footprint is
a different transformation from enlarging one cone. The proposed drop tests
must check both the terminal regime and orientation.
\item A shrinking-suffix bubble-sort implementation in Problem 4.25 uses
at most \(n-1\) passes and \(n(n-1)/2\) comparisons in its fixed-pass worst
case. Problem 4.26 starts with
\((98,33,34,62,31,58,61,15)\); the two sorted halves are
\((33,34,62,98)\) and \((15,31,58,61)\), yielding
\((15,31,33,34,58,61,62,98)\). The recurrence below counts the merge work.
\item Problem 4.31's transformation \(v\mapsto\sqrt{A/B}/v\) is reflection
of \(\log v\) about \(\log v_E\). Problem 4.33 defines the lift coefficient
using wing area: under steady vertical balance,
\(c_L=2mg/(\rho v^2A_{\rm wing})\). It is not the coefficient based on
cross-sectional area.
\item Problem 4.37 supplies a \(0.4\)-kilogram bird and a source-reported
journey of 11680 kilometres in 8.1 days, whose average ground speed is about
\(16.7\,\mathrm{m/s}\). Ground speed and minimum-energy airspeed need not
coincide. Problem 4.40 supplies a 12-milligram housefly; wing scale,
air density and speed remain additional inputs to the proposed power
comparison. Problem 4.38 requires the full resistance ratio
\((d_1/d_2)(\kappa_2/\kappa_1)(A_2/A_1)\), including conductivities
\(\kappa_1,\kappa_2\), not thickness and area alone.
\item In Problem 5.11, the selected deep-water variables
\((v,g,\omega,\rho)\) allow \(v\omega/g\), with density excluded because
it alone carries mass. The additional supplied coefficient one gives
\(v=gT/(2\pi)\), about \(26.5\,\mathrm{m/s}\) for a 17-second period
and \(g=9.8\,\mathrm{m/s^2}\). Equality to a generating wind speed is
not a dimensional conclusion. In Problem 5.15, increasing shallow-water
depth from 1 to 4 centimetres doubles the predicted speed if its coefficient
and long-wave regime remain fixed; ocean travel time further requires a
path length and a depth model.
\item A thin ring in Problem 5.18 has \(I/(mr^2)=1\), giving
\(a=g\sin\theta/2\) in the no-slip model, independently of its mass or
size. In Problem 5.20, \(GM/(rc^2)\) compares the Newtonian potential-energy
scale \(GMm/r\) with \(mc^2\). Problem 5.22 requests the corresponding
coefficient in Brans--Dicke theory; that is an additional physical-theory
question, not another consequence of the dimension matrix.
\item For Problem 5.23, a stipulated speed \(1\,\mathrm{m/s}\) and
viscosity \(1.5\,10^{-5}\,\mathrm{m^2/s}\) require diameter
\(4.5\,\mathrm m\) for Reynolds number \(3\,10^5\). Keeping the speed
fixed at that size is a further assumption. Problem 5.24 asks for the
corresponding estimates for a raindrop, mosquito, walking person and
cruising jet; each requires its own speed, length convention and fluid.
\end{itemize}


\begin{hbtheorem}{thm:p15-dimensional-source-extensions}{proved=annex:p15.dimensional.source-extensions}{Selected source extensions under explicit models}
\begin{enumerate}
\item For \(n=2^h\), a comparison-cost model
\(C(n)=2C(n/2)+(n-1)\), \(C(1)=0\), gives
\(C(n)=n\log_2n-n+1\). A product-count model
\(M(n)=3M(n/2)\), \(M(1)=1\), gives
\(M(n)=n^{\log_2 3}\); the corresponding four-product model gives \(n^2\).
\item Under a uniform-sphere enclosed-mass law and inverse-square gravity,
the interior field is proportional to radius and the exterior field to
inverse radius squared.
\item For a nonzero oscillation governed by \(m\ddot x=-4Cx^3\),
\(C,m>0\), with turning amplitude \(x_0>0\), the period is
\[
 T=4\sqrt{\frac m{2C}}\,x_0^{-1}
       \int_0^1\frac{du}{\sqrt{1-u^4}}.
\]
\item Under translational force balance \(mg\sin\theta-f=ma\),
torque balance \(fr=I a/r\), and no slip, the acceleration is
\(a=g\sin\theta/(1+I/(mr^2))\), for \(m,r>0\), \(I\geq0\).
\end{enumerate}
\end{hbtheorem}

These recurrence statements specify abstract cost models. They do not assert
that every implementation has exactly those counts or that integer
multiplication ignores carry handling. The physical extensions likewise
give consequences of their stated equations.

\section{Using the method as a reference}
\label{sec:p15-dimensional-reference}

Begin with the target relation and its domain. Check units and the relevance
of the comparison; list omitted constants and dimensionless controls.
Compute the rank and products, choose a convenient output representation,
then identify the additional law, calibration or regime information that
determines the remaining function. In a scaling comparison state what remains
fixed. In an estimate state its input ranges and distinguish those ranges
from probability intervals.

The measurement-uncertainty treatment
(\cref{sec:p15-measurement-uncertainty}) supplies a different calculation:
propagating a stated uncertainty model through a function. Dimensional
consistency alone supplies neither that probability model nor its covariance.
Approximation and limiting-case methods are developed in
\cref{sec:p15-approximate-quantitative-inference}; they cannot turn a guessed
dimensionless constant into an exact result.
